\documentclass[11pt]{article}

% ============================================================
% Packages
% ============================================================
\usepackage[utf8]{inputenc}
\usepackage[margin=1in]{geometry}
\usepackage{amsmath,amssymb,amsthm}
\usepackage{graphicx}
\usepackage{booktabs}
\usepackage{xcolor}
\usepackage{hyperref}
\usepackage{tikz}
\usepackage{enumitem}
\usepackage{pgfplots}

\pgfplotsset{compat=1.18}

% ============================================================
% Document Formatting
% ============================================================
\setlength{\parskip}{0.6em}
\setlength{\parindent}{0pt}

% ============================================================
% List Formatting
% ============================================================
\setlist[enumerate]{
    leftmargin=2em,
    itemsep=0.25em
}

% ============================================================
% Solution Environment
% ============================================================
\newenvironment{solution}
    {\par\noindent\textbf{Solution.}\par}
    {\par}

% ============================================================
% Link Styling
% ============================================================
\hypersetup{
    colorlinks=true,
    linkcolor=blue,
    urlcolor=cyan
}

% ============================================================
% Metadata
% ============================================================
\title{Ordinary Differential Equations: Abridged Notes}
\author{Chris Ritter}
\date{September 10, 2026}

% ============================================================
% Document
% ============================================================

\begin{document}
\maketitle
\begin{center}
\small\emph{These abridged notes were reconstructed afterward because I did not take
direct notes during class. They are not a record of the actual lecture.}
\end{center}
\tableofcontents
\newpage

\section{Topic Summary}
An autonomous equation has the form \(y'=f(y)\): its rate depends on
the current state, not explicitly on time. Its equilibria are the
roots of \(f\). A phase line records the sign of \(f\) between
consecutive equilibria and shows whether nearby solutions move toward
or away from each one. This gives qualitative information about
population models without solving the equation.

\section{Phase Lines and Stability}
To analyze \(y'=f(y)\), first solve \(f(y)=0\). On each interval
between equilibrium values, test one point to find the sign of \(f\):
positive means solutions increase; negative means they decrease.
An equilibrium is \textbf{asymptotically stable} if nearby solutions
approach it as time increases, and \textbf{unstable} if nearby
solutions move away. When \(f'(y_*)\ne0\), the sign of \(f'(y_*)\)
gives the same local test: negative is stable and positive is unstable.

\subsection{Example: Quadratic Population Rate}
\textbf{Adapted from Boyce et al., 11th ed., \S2.5, Problem 1.}
Consider \(y'=ay+by^2\) with \(a,b>0\).
\begin{solution}
Factor \(f(y)=y(a+by)\). The equilibria are \(y=-a/b\) and \(y=0\).
The sign chart is
\[
\begin{array}{c|ccc}
y&(-\infty,-a/b)&(-a/b,0)&(0,\infty)\\\hline
f(y)&+&-&+
\end{array}
\]
Thus arrows point toward \(-a/b\) from both sides, while they point
away from 0. Therefore
\[
\boxed{-a/b\text{ is asymptotically stable, and }0\text{ is unstable}.}
\]
For nonnegative population, \(y=0\) is the only equilibrium in the
physical range; every positive solution initially increases. The
negative equilibrium is mathematically valid but not a population.
\end{solution}

\subsection{Example: Three Equilibria}
\textbf{Adapted from Boyce et al., 11th ed., \S2.5, Problem 2.}
Consider \(y'=y(y-1)(y-2)\) for \(y\geq0\).
\begin{solution}
The equilibria are 0, 1, and 2. Testing one value in each interval
gives
\[
\begin{array}{c|cccc}
y&(-\infty,0)&(0,1)&(1,2)&(2,\infty)\\\hline
f(y)&-&+&-&+
\end{array}
\]
The arrows on \((0,1)\) point right and those on \((1,2)\) point
left, so nearby solutions approach 1. On either side of 0 and 2,
nearby solutions move away. Hence
\[
\boxed{1\text{ is asymptotically stable; }0\text{ and }2
\text{ are unstable}.}
\]
For physical initial data \(y_0\geq0\), the negative part of the
phase line is only a mathematical extension.
\end{solution}
\end{document}
